\section{Finite Energy-Shell Projection and Effective Hamiltonians}
\label{sec:ch21-finite-effective}

The preceding sections built finite operators before asking which states they
preserve.  In particular, the Umklapp word has a definite change of four
species charges, but this selection rule is not a spectral statement, and its
nonzero action need not preserve a chosen finite-energy carrier.  We now face
a different, constructive question: \emph{if we intentionally retain only a
low-energy part of a finite Hilbert space, what operator describes its
interaction with the discarded states?}  Simply deleting the discarded
states removes the very virtual excursions that can modify low-energy
matrix elements.  A finite Schrieffer--Wolff transformation supplies a
controlled algebraic answer through second order in a coupling parameter.

This section develops the method rather than presupposing familiarity with
renormalization.  The starting objects are ordinary finite-dimensional
Hilbert spaces, orthogonal projections and matrix commutators.  We shall
separate (i) exact partition identities, (ii) a solvable non-resonant matrix
equation, (iii) an \emph{exact identity between polynomials of degree two},
and (iv) approximation of an analytic unitary rotation.  Their hypotheses
are different.  In particular, none of these steps by itself constructs a
charge gap or a continuum renormalization-group (RG) flow.

\subsection{What does an energy shell discard?}
\label{subsec:ch21-energy-shell}

Fix a finite lattice of length $L=2h$ with $h\in\mathbb N_{>0}$, a
positive-mode cutoff $M\geq1$, a set of species labels (in this chapter,
$\{R,L\}\times\{\uparrow,\downarrow\}$), and an admissible integer
species-charge vector $\boldsymbol N$.  The finite CAR Fock carrier will be
denoted $\mathcal F_L$.  We use the occupation-space excitation grading
$\mathsf E\colon\{\text{occupation configurations}\}\to\mathbb N_0$,
relative to the specified charge-sector ground configuration.  The
charge-sector \emph{budget} is the span
\begin{equation}
\mathcal B_{\boldsymbol N,K}
 :=\operatorname{span}\bigl\{\lvert A\rangle:
 \boldsymbol N(A)=\boldsymbol N,\ \mathsf E(A)\leq K\bigr\}
 \subseteq\mathcal F_L,
\label{eq:ch21-budget-definition}
\end{equation}
where $K\in\mathbb N_0$.  We emphasize that \eqref{eq:ch21-budget-definition}
is a span of occupation kets selected by a scalar \emph{grade}; it is not
a vector inequality of the form $H\psi\leq K\psi$.  The allowed grades,
charges and finite-volume margins are those of the earlier energy-budget
construction (see Chapter~\ref{ch:vacuum_budget}).

For $K\geq1$, let $\Pi_{\boldsymbol N,K}$ be the orthogonal coordinate
projection onto this budget and define the two \emph{ambient} projections
\begin{align}
P&:=\Pi_{\boldsymbol N,K-1},\\
Q&:=\Pi_{\boldsymbol N,K}-\Pi_{\boldsymbol N,K-1}
  =\operatorname{proj}_{\mathcal H_{\boldsymbol N,K}},\qquad
\mathcal H_{\boldsymbol N,K}:=
 \operatorname{span}\{\lvert A\rangle:\boldsymbol N(A)=\boldsymbol N,
                                  \mathsf E(A)=K\}.
\label{eq:ch21-shell-projections}
\end{align}
Both projections act on $\mathcal F_L$, even though their ranges are
subspaces of $\mathcal B_{\boldsymbol N,K}$.  Their defining coordinate
supports are disjoint.

\begin{proposition}[Exact finite shell partition]
\label{prop:ch21-shell-partition}
For $K\geq1$ the ambient projections satisfy
\begin{equation}
 P^2=P=P^\dagger,\quad Q^2=Q=Q^\dagger,
 \quad PQ=QP=0,\quad P+Q=\Pi_{\boldsymbol N,K}.
\label{eq:ch21-ambient-partition}
\end{equation}
Consequently, with $\mathcal H:=\mathcal B_{\boldsymbol N,K}$ as the
\emph{new carrier}, the restricted operators obey $P+Q=I_{\mathcal H}$
and give the orthogonal decomposition
$\mathcal H=\operatorname{ran}P\oplus\operatorname{ran}Q$.
\end{proposition}
\begin{proof}
In the orthonormal occupation basis, every coordinate projector acts as
identity on a selected ket and as zero on the others.  Its square therefore
equals itself, and its matrix is real diagonal, hence self-adjoint.  The
supports $\mathsf E\leq K-1$ and $\mathsf E=K$ are disjoint: applying
$PQ$ or $QP$ to each basis ket gives zero.  Their union is exactly the
support $\mathsf E\leq K$, within the fixed charge sector.  Thus their
sum equals $\Pi_{\boldsymbol N,K}$ on every occupation ket, proving
\eqref{eq:ch21-ambient-partition}.  This last projection is the identity
\emph{only after} restricting the carrier to its range.  Orthogonality and
direct-sum decomposition follow by writing every $\psi\in\mathcal H$ as
$\psi=P\psi+Q\psi$ and observing that the ranges intersect trivially.
\end{proof}

\begin{example}[Removing energy two is not removing mode two]
\label{ex:ch21-shell-versus-mode}
Consider the polynomial oscillator carrier with two positive modes,
$X_1$ and $X_2$, weighted by energies $1$ and $2$, respectively.  For
$M\geq2$, the illustrative filtered monomial basis at budget two is
\begin{equation}
 \mathcal B_2^{\mathrm{osc}}
  =\operatorname{span}\{1,X_1,X_1^2,X_2\},\qquad
 \mathcal B_1^{\mathrm{osc}}=\operatorname{span}\{1,X_1\}.
\label{eq:ch21-osc-shell-example}
\end{equation}
The shell removed by $P$ is
$\operatorname{span}\{X_1^2,X_2\}$.  Deleting the \emph{mode} $X_2$
instead would retain $X_1^2$ and yield
$\operatorname{span}\{1,X_1,X_1^2\}$.  These are distinct projections.
Their difference persists in a concrete finite-Fock realization: under
the charge-admissibility and frozen-margin hypotheses required for the
repeated positive mode-one density excitation, the vector
$\rho_{1,\nu}^{\,2}\lvert\boldsymbol N\rangle_0$ is nonzero, has
excitation grade two, is killed by the budget-one projector, and is
fixed by the exact-grade-two shell projector.  Here $\rho_{1,\nu}$
uses the convention that positive $m$ raises the excitation grade.
The needed nonvanishing is a finite CAR/current statement, not an
assumption that the polynomial oscillator is literally the full Fock
carrier.  The supporting Core A10 result gives this witness with
$h\geq2$ and its stated admissibility/margin conditions.
\end{example}

The letter $K$ will continue to denote the \emph{energy budget}, not a
Luttinger parameter.  The shell splitting does not assert that its
grade is the interacting Hamiltonian's spectrum: the grade is a basis
filtration fixed in advance.  A Hamiltonian may mix these spaces, which
is precisely why an effective Hamiltonian is needed.

\subsection{Block matrices and the denominator hypothesis}
\label{subsec:ch21-block-setup}

Work now on an arbitrary finite-dimensional complex Hilbert space
$\mathcal H=\mathcal P\oplus\mathcal Q$, with orthogonal projections
$P$ and $Q=I_{\mathcal H}-P$.  This could be the restricted budget
carrier from Proposition~\ref{prop:ch21-shell-partition}; alternatively it
could be a different explicitly selected finite spectral window.  Let
\begin{equation}
H(t)=H_0+tV,\qquad H_0^\dagger=H_0,\quad V^\dagger=V,
\quad t\in\mathbb R,\quad [P,H_0]=0.
\label{eq:ch21-H-perturbation}
\end{equation}
The last hypothesis means that $H_0$ preserves $\mathcal P$ and
$\mathcal Q$.  It is not automatic for a chosen interacting baseline.
In particular, an interbranch pairing term can mix bare Haldane
partitions: one must diagonalize the proposed $H_0$ on an appropriate
carrier or use a baseline, such as the free Sugawara Hamiltonian, for
which the requisite invariance actually holds.

Insert $I=P+Q$ to write each operator in a $2\times2$ array of blocks:
\begin{equation}
H_0=\begin{pmatrix}A&0\\0&D\end{pmatrix},\qquad
V=\begin{pmatrix}V_{PP}&B\\B^\dagger&V_{QQ}\end{pmatrix},
\quad A=PH_0P\big|_{\mathcal P},\quad D=QH_0Q\big|_{\mathcal Q},
\label{eq:ch21-block-matrices}
\end{equation}
where $B=PVQ\colon\mathcal Q\to\mathcal P$.  Define
\begin{equation}
 V_{\mathrm d}:=PVP+QVQ,\qquad
 V_{\mathrm o}:=PVQ+QVP,\qquad V=V_{\mathrm d}+V_{\mathrm o}.
\label{eq:ch21-diag-offdiag}
\end{equation}
The subscripts refer to this \emph{chosen decomposition}, not to
momentum conservation or locality.  A block-off-diagonal perturbation
moves a state between retained and discarded sectors.  Its second-order
virtual excursion may start and end in $\mathcal P$ even when the first
order transition does not stay there.

Because $H_0$ is self-adjoint and preserves the two summands, we may
choose an orthonormal eigenbasis
$\{\lvert\lambda\rangle\}_{\lambda\in I_P}$ of $A$ and
$\{\lvert\mu\rangle\}_{\mu\in I_Q}$ of $D$,
with \emph{real} eigenvalues $E_\lambda$ and $E_\mu$.  No relation between
the two eigenvalue sets has yet been assumed.  A virtual transition
requires an energy denominator; it is therefore necessary to state the
non-resonance condition rather than dividing by a possibly zero number.

\begin{definition}[Coupling-wise non-resonance]
\label{def:ch21-coupling-nonresonance}
The pair $(H_0,V)$ is \emph{non-resonant across $P/Q$ on its nonzero
couplings} if
\begin{equation}
 \langle\lambda\lvert V\rvert\mu\rangle\ne0
 \quad\Longrightarrow\quad E_\lambda\ne E_\mu,
 \qquad \lambda\in I_P,\ \mu\in I_Q.
\label{eq:ch21-nonresonance}
\end{equation}
Degenerate pairs with zero cross-block matrix element are allowed.
\end{definition}

The stronger assumption
$\operatorname{spec}(A)\cap\operatorname{spec}(D)=\varnothing$
guarantees invertibility of the Sylvester operator introduced below.
The weaker, coupling-wise assumption is sufficient for \emph{existence}
of the particular solution used in this section, but not generally for
uniqueness.  Since there are finitely many coupled pairs, its nonzero
energy differences have a strictly positive minimum when any such
pairs exist; that finite number need not remain bounded away from zero
as $L$ or the budget varies.

\subsection{Solving the Sylvester equation}
\label{subsec:ch21-sylvester}

We seek a first-order change of basis generated by an anti-Hermitian,
block-off-diagonal operator $S_1$:
\begin{equation}
 S_1^\dagger=-S_1,\qquad PS_1P=QS_1Q=0,
 \qquad [S_1,H_0]=-V_{\mathrm o}.
\label{eq:ch21-generator-equation}
\end{equation}
Anti-Hermiticity ensures $e^{tS_1}$ is unitary for real $t$, while the
commutator condition cancels the first-order coupling between
$\mathcal P$ and $\mathcal Q$.  It is important that the commutator is
ordered $[S_1,H_0]=S_1H_0-H_0S_1$; reversing it reverses the
energy-denominator sign.

Write
$S_1=\bigl(\begin{smallmatrix}0&X\\-X^\dagger&0\end{smallmatrix}\bigr)$,
where $X:\mathcal Q\to\mathcal P$.  Multiplying the block matrices in
\eqref{eq:ch21-block-matrices} gives
\begin{equation}
 [S_1,H_0]_{PQ}=XD-AX=-B
 \quad\Longleftrightarrow\quad AX-XD=B.
\label{eq:ch21-sylvester}
\end{equation}
This is a \emph{Sylvester equation}: a linear equation for an unknown
operator $X$ in which fixed operators multiply it on the left and right.
It is not a division by the matrix $A-D$, since the matrices can act on
different spaces and need not commute.

\begin{theorem}[Finite non-resonant Schrieffer--Wolff generator]
\label{thm:ch21-sw-generator}
Under \eqref{eq:ch21-H-perturbation} and
Definition~\ref{def:ch21-coupling-nonresonance}, define the
cross-block matrix entries
\begin{equation}
\langle\lambda\lvert X\rvert\mu\rangle:=
\begin{cases}
\displaystyle\frac{\langle\lambda\lvert V\rvert\mu\rangle}
 {E_\lambda-E_\mu}, &
       \langle\lambda\lvert V\rvert\mu\rangle\ne0,\\[6pt]
0,&\langle\lambda\lvert V\rvert\mu\rangle=0,
\end{cases}
\qquad \lambda\in I_P,\ \mu\in I_Q.
\label{eq:ch21-X-entries}
\end{equation}
Then $S_1=\bigl(\begin{smallmatrix}0&X\\-X^\dagger&0\end{smallmatrix}\bigr)$
is anti-Hermitian, block-off-diagonal, and solves
$[S_1,H_0]=-V_{\mathrm o}$.  The zero choice on uncoupled pairs remains
well defined even when their energies coincide.
\end{theorem}
\begin{proof}
The denominator is nonzero in the first case by
\eqref{eq:ch21-nonresonance}; the second case involves no division.
In the chosen eigenbases,
\begin{align*}
\langle\lambda\lvert AX-XD\rvert\mu\rangle
 &=(E_\lambda-E_\mu)
     \langle\lambda\lvert X\rvert\mu\rangle\\
 &=\langle\lambda\lvert V\rvert\mu\rangle,
\end{align*}
including when the matrix element of $V$ is zero.  Thus $AX-XD=B$
entrywise.  The opposite off-diagonal block of
$[S_1,H_0]$ is the adjoint of the first block, because
$S_1^\dagger=-S_1$ and $H_0^\dagger=H_0$ make $[S_1,H_0]$ Hermitian.
More explicitly, the opposite entry is
\begin{equation*}
 \langle\mu\lvert S_1\rvert\lambda\rangle
 =-\overline{\langle\lambda\lvert S_1\rvert\mu\rangle}
 =\frac{\langle\mu\lvert V\rvert\lambda\rangle}{E_\mu-E_\lambda}
\end{equation*}
when coupled, and zero otherwise.  Both diagonal blocks of
$[S_1,H_0]$ vanish because $H_0$ is block diagonal and $S_1$ is
off diagonal.  Collecting the four blocks proves the commutator
equation; the displayed block form proves anti-Hermiticity directly.
\end{proof}

\begin{remark}[Why the solution need not be unique]
If $E_\lambda=E_\mu$ for an uncoupled cross-block pair, the
corresponding entry of the homogeneous equation
$AX-XD=0$ is unconstrained.  Setting it to zero is a convenient,
explicit choice, not a uniqueness theorem.  If the spectra of $A$ and
$D$ are disjoint, the homogeneous equation has only the zero solution,
and the constructed $X$ is unique.
\end{remark}

\subsection{The ordered second-order transformation}
\label{subsec:ch21-sw-bch}

For any finite matrices $S$ and $T$, write
$\operatorname{ad}_S(T):=[S,T]$ and
$\operatorname{ad}_S^j(T)$ for its $j$-fold iteration.  We rotate the
Hamiltonian according to
\begin{equation}
 \widetilde H(t):=e^{tS_1}(H_0+tV)e^{-tS_1}.
\label{eq:ch21-sw-rotated}
\end{equation}
The orientation is fixed: we conjugate by $e^{+tS_1}$ on the left.
The finite-matrix exponential exists for every real $t$ and the
rotated Hamiltonian is self-adjoint and isospectral to $H(t)$, even
though the compression $P\widetilde H(t)P$ need not reproduce the exact
low-energy spectrum at a finite, non-small $t$.

To obtain the perturbative expansion, differentiate
$e^{tS_1}Te^{-tS_1}$ at $t=0$.  Its successive derivatives are
$\operatorname{ad}_{S_1}^j(T)$.  Hence the ordered polynomial through
degree two is
\begin{align}
B_2(t):={}&H_0+t\bigl(V+[S_1,H_0]\bigr)\notag\\
 &+t^2\left([S_1,V]
        +\frac12[S_1,[S_1,H_0]]\right).
\label{eq:ch21-bch-two-polynomial}
\end{align}
The equality between $B_2$ and the Taylor polynomial is a finite
algebraic fact, irrespective of whether $t$ is small.  The claim that
it approximates $\widetilde H(t)$ needs a separate remainder statement.

\begin{theorem}[Exact compressed degree-two identity]
\label{thm:ch21-sw-compressed}
Assume that $P,Q$ are complementary orthogonal projections on
$\mathcal H$, $H_0$ commutes with $P$, and $S_1$ obeys
\eqref{eq:ch21-generator-equation}.  Then, as an \emph{exact polynomial
identity} on $\mathcal P=\operatorname{ran}P$,
\begin{equation}
 P B_2(t)P
 =PH_0P+tPVP+\frac{t^2}{2}P[S_1,V]P.
\label{eq:ch21-sw-polynomial}
\end{equation}
The right-hand side is the second-order Schrieffer--Wolff effective
Hamiltonian, denoted $H_{\mathrm{eff}}^{[2]}(t)$.  No exponential
appears in the asserted equality.
\end{theorem}
\begin{proof}
Substitute $[S_1,H_0]=-V_{\mathrm o}$ into
\eqref{eq:ch21-bch-two-polynomial}.  The first-order coefficient is
$V-V_{\mathrm o}=V_{\mathrm d}$; in particular,
$PV_{\mathrm d}P=PVP$.  For the second-order term,
\begin{align*}
[S_1,V]+\tfrac12[S_1,[S_1,H_0]]
 &= [S_1,V]-\tfrac12[S_1,V_{\mathrm o}]\\
 &= [S_1,V_{\mathrm d}]+\tfrac12[S_1,V_{\mathrm o}].
\end{align*}
A product of an off-diagonal and a diagonal block operator is
off-diagonal, so $P[S_1,V_{\mathrm d}]P=0$.  Therefore the compressed
coefficient becomes
$\tfrac12P[S_1,V_{\mathrm o}]P$.  Moreover,
$P[S_1,V]P=P[S_1,V_{\mathrm o}]P$ by the same block argument.
The constant term is $PH_0P$, yielding exactly
\eqref{eq:ch21-sw-polynomial}.  Every equality used here is an
ordinary finite-matrix identity, with no perturbative limiting
argument.
\end{proof}

\begin{proposition}[Intermediate-state formula]
\label{prop:ch21-sw-intermediate}
For $\lambda,\lambda'\in I_P$ the second-order correction has matrix
entries
\begin{equation}
\begin{aligned}
\langle\lambda\lvert W^{(2)}\rvert\lambda'\rangle
  &=\frac12\sum_{\substack{\mu\in I_Q\\
       V_{\lambda\mu}V_{\mu\lambda'}\ne0}}
  V_{\lambda\mu}V_{\mu\lambda'}
  \left(\frac{1}{E_\lambda-E_\mu}
        +\frac{1}{E_{\lambda'}-E_\mu}\right),\\
W^{(2)}&:=\frac12P[S_1,V]P,
\qquad V_{ab}:=\langle a\lvert V\rvert b\rangle.
\end{aligned}
\label{eq:ch21-sw-intermediate-sum}
\end{equation}
Only nonzero coupled terms are included in the displayed sum, so its
denominators are defined under \eqref{eq:ch21-nonresonance}.  In
particular,
\begin{equation}
\langle\lambda\lvert W^{(2)}\rvert\lambda\rangle
 =\sum_{\substack{\mu\in I_Q\\V_{\lambda\mu}\ne0}}
 \frac{|V_{\lambda\mu}|^2}{E_\lambda-E_\mu}.
\label{eq:ch21-sw-diagonal-shift}
\end{equation}
\end{proposition}
\begin{proof}
Since $S_1$ is off diagonal, the $PP$ block of $S_1V$ is
$S_{1,PQ}V_{QP}$ and that of $VS_1$ is $V_{PQ}S_{1,QP}$.
Insert the orthonormal resolution $I_{\mathcal Q}=\sum_{\mu\in I_Q}
\lvert\mu\rangle\langle\mu\rvert$ to get
\begin{equation*}
\langle\lambda\lvert[S_1,V]\rvert\lambda'\rangle
=\sum_{\mu\in I_Q}\bigl((S_1)_{\lambda\mu}
 V_{\mu\lambda'}-V_{\lambda\mu}(S_1)_{\mu\lambda'}\bigr).
\end{equation*}
Theorem~\ref{thm:ch21-sw-generator} gives
$(S_1)_{\lambda\mu}=V_{\lambda\mu}/(E_\lambda-E_\mu)$
on coupled pairs and
$(S_1)_{\mu\lambda'}=-V_{\mu\lambda'}/(E_{\lambda'}-E_\mu)$.
Terms with vanishing couplings contribute zero by definition.  This
yields \eqref{eq:ch21-sw-intermediate-sum}.  Taking
$\lambda'=\lambda$ and using Hermiticity of $V$ proves
\eqref{eq:ch21-sw-diagonal-shift}.
\end{proof}

The denominator in \eqref{eq:ch21-sw-diagonal-shift} explains the
physical interpretation: the retained state makes a virtual excursion
to a discarded eigenstate and returns.  In a strictly ordered baseline
spectrum $E_\lambda<E_\mu$ for every coupled pair, each diagonal
second-order shift is nonpositive.  Without that ordering, cancellations
are possible.  Thus even the \emph{sign} of the correction is not
determined by the mere existence of a $P/Q$ splitting.

\subsection{What is the remainder, exactly?}
\label{subsec:ch21-sw-remainder}

There are two different uses of the notation $O(t^3)$.  In a formal
power series ring $\operatorname{End}(\mathcal H)[[t]]$, the statement
\begin{equation}
P\widetilde H(t)P\equiv H_{\mathrm{eff}}^{[2]}(t)
 \pmod{t^3}
\label{eq:ch21-sw-formal-mod}
\end{equation}
means \emph{coefficientwise} equality in degrees zero, one and two.
It does not evaluate an operator at a numerical $t$, and it does not
by itself bound the discarded terms.  For the finite analytic
exponential we can additionally obtain an explicit bound; this is
an elementary finite-matrix consequence, stated here separately from
the supplied A10 Lean theorem.

\begin{proposition}[Finite analytic Taylor remainder]
\label{prop:ch21-sw-analytic-remainder}
Under the hypotheses of Theorem~\ref{thm:ch21-sw-compressed}, for real
$t$ let
\begin{equation}
 R_3(t):=\widetilde H(t)-B_2(t).
\label{eq:ch21-rem-definition}
\end{equation}
Then the exact compressed analytic Hamiltonian obeys
\begin{equation}
 P\widetilde H(t)P=H_{\mathrm{eff}}^{[2]}(t)+PR_3(t)P,
\label{eq:ch21-sw-exact-plus-remainder}
\end{equation}
where, writing $\mathcal U_r(T)=e^{rS_1}Te^{-rS_1}$,
\begin{align}
R_3(t)={}&\frac{t^3}{2}\int_0^1(1-u)^2
  \mathcal U_{ut}\bigl(\operatorname{ad}_{S_1}^3(H_0)\bigr)\,du\notag\\
 &+t^3\int_0^1(1-u)
  \mathcal U_{ut}\bigl(\operatorname{ad}_{S_1}^2(V)\bigr)\,du.
\label{eq:ch21-sw-integral-remainder}
\end{align}
For the operator norm, using $S_1^\dagger=-S_1$,
\begin{equation}
\|PR_3(t)P\|\leq\|R_3(t)\|
\leq |t|^3\left(
\frac{\|\operatorname{ad}_{S_1}^3(H_0)\|}{6}
+\frac{\|\operatorname{ad}_{S_1}^2(V)\|}{2}\right).
\label{eq:ch21-sw-norm-bound}
\end{equation}
The bound is finite-volume and its constant depends on the chosen
carrier, $H_0,V$ and the selected generator $S_1$.
\end{proposition}
\begin{proof}
For a fixed $T$, put $F_T(t)=e^{tS_1}Te^{-tS_1}$.
Differentiating the matrix exponential gives
$F_T^{(j)}(r)=\mathcal U_r(\operatorname{ad}_{S_1}^jT)$.
Taylor's theorem with integral remainder through order two for
$F_{H_0}$ gives
\begin{equation*}
F_{H_0}(t)=H_0+t\operatorname{ad}_{S_1}(H_0)
+\frac{t^2}{2}\operatorname{ad}_{S_1}^2(H_0)
+\frac{t^3}{2}\int_0^1(1-u)^2
  \mathcal U_{ut}(\operatorname{ad}_{S_1}^3H_0)\,du.
\end{equation*}
Taylor's theorem through order one for $F_V$ gives
\begin{equation*}
F_V(t)=V+t\operatorname{ad}_{S_1}(V)
+t^2\int_0^1(1-u)
 \mathcal U_{ut}(\operatorname{ad}_{S_1}^2V)\,du.
\end{equation*}
Add $F_{H_0}(t)+tF_V(t)$ and compare with
\eqref{eq:ch21-bch-two-polynomial}; this proves
\eqref{eq:ch21-sw-integral-remainder} and then
\eqref{eq:ch21-sw-exact-plus-remainder}.
For real $t$ and $u$, $e^{utS_1}$ is unitary, so
$\|\mathcal U_{ut}(T)\|=\|T\|$.  Apply the triangle inequality and
$\int_0^1(1-u)^2du=1/3$ and
$\int_0^1(1-u)du=1/2$ to obtain
\eqref{eq:ch21-sw-norm-bound}.  Finally, orthogonal projections are
contractions, hence $\|PR_3P\|\leq\|R_3\|$.
\end{proof}

\begin{remark}[A finite error is not a uniform scaling estimate]
The right side of \eqref{eq:ch21-sw-norm-bound} may increase with
lattice size or as the non-resonant denominators become small.  It is
a genuine estimate at fixed finite carrier, but not a proof of an
$L\to\infty$ limit or a scale-independent perturbative expansion.
The coefficientwise polynomial result of the supplied Core A10 module
should not be reported as formal verification of this analytic bound.
\end{remark}

\subsection{A two-level calculation, including what fails at finite coupling}
\label{subsec:ch21-two-level}

\begin{example}[Explicit Schrieffer--Wolff rotation]
\label{ex:ch21-sw-two-level}
Let $\mathcal H=\mathbb C^2$, $P=\operatorname{diag}(1,0)$ and
$Q=\operatorname{diag}(0,1)$.  Take
\begin{equation}
 H_0=\begin{pmatrix}0&0\\0&1\end{pmatrix},\qquad
 V=\begin{pmatrix}0&1\\1&0\end{pmatrix}.
\label{eq:ch21-sw-2x2-input}
\end{equation}
The low state has energy zero and the high state energy one.  The
Sylvester denominator is $0-1=-1$, so
\begin{equation}
 S_1=\begin{pmatrix}0&-1\\1&0\end{pmatrix},\qquad
 [S_1,H_0]=-V,
 \quad[S_1,V]=\begin{pmatrix}-2&0\\0&2\end{pmatrix}.
\label{eq:ch21-sw-2x2-generator}
\end{equation}
The effective Hamiltonian is therefore the scalar operator
\begin{equation}
 H_{\mathrm{eff}}^{[2]}(t)=-t^2 P
 \quad\text{on }\operatorname{ran}P.
\label{eq:ch21-sw-2x2-effective}
\end{equation}
We can test exactly which part of this assertion is an approximation.
Since $S_1^2=-I$, the rotation is
\begin{equation}
 e^{tS_1}=\begin{pmatrix}\cos t&-\sin t\\
                                   \sin t&\cos t\end{pmatrix}.
\end{equation}
Direct multiplication gives the \emph{exact} compressed entry
\begin{align}
 (P\widetilde H(t)P)_{11}
 &=\sin^2t-t\sin(2t)
   =-t^2+t^4+O(t^6),
\label{eq:ch21-sw-2x2-rotated}\end{align}
and the off-diagonal entry
\begin{equation}
 (\widetilde H(t))_{12}
 =t\cos(2t)-\tfrac12\sin(2t)
 =-\tfrac43t^3+O(t^5).
\label{eq:ch21-sw-2x2-off}
\end{equation}
Thus the first-order mixing is canceled but the chosen rotation is
not an exact block diagonalization for arbitrary $t$.  Independently,
the exact eigenvalues of $H_0+tV$ are
\begin{equation}
 E_\pm(t)=\frac{1\pm\sqrt{1+4t^2}}2.
\label{eq:ch21-sw-2x2-eigvals}
\end{equation}
At $t=2/3$, the lower eigenvalue is $-1/3$ (with eigenvector
$(-2,1)$), whereas \eqref{eq:ch21-sw-2x2-effective} gives $-4/9$.
The two values differ because a second-order effective operator is
not an exact finite-coupling spectral identification.
\end{example}

The same caution applies when the quadratic correction happens to
vanish.  A nonzero perturbation can couple retained and discarded
states through \emph{different} intermediate energies with cancelling
denominators.

\begin{example}[Zero second-order correction with nonzero leakage]
\label{ex:ch21-sw-three-level}
Take $\mathcal P=\operatorname{span}\{e_1\}$,
$\mathcal Q=\operatorname{span}\{e_2,e_3\}$ and
\begin{equation}
 H_0=\operatorname{diag}(0,-1,1),\quad
 V=\begin{pmatrix}0&1&1\\1&0&0\\1&0&0\end{pmatrix},\quad
 S_1=\begin{pmatrix}0&1&-1\\-1&0&0\\1&0&0\end{pmatrix}.
\label{eq:ch21-sw-three-input}
\end{equation}
Here $PVQ\ne0$ and $[S_1,H_0]=-V$, but
\begin{equation}
\langle e_1\lvert W^{(2)}\rvert e_1\rangle
 =\frac{|V_{12}|^2}{0-(-1)}
  +\frac{|V_{13}|^2}{0-1}=1-1=0.
\label{eq:ch21-sw-three-cancel}
\end{equation}
A vanishing compressed second-order commutator therefore does
\emph{not} imply $PVQ=QVP=0$.  Conversely, if both off-block
couplings vanish, choosing $S_1=0$ gives the exact identity-rotation
compression $P(H_0+tV)P=PH_0P+tPVP$ and zero correction.  That
forward implication does not require a non-resonance hypothesis.
\end{example}

\subsection{Scope of the finite construction and the remaining chapter obligations}
\label{subsec:ch21-effective-scope}

The discrete procedure has three mathematically independent inputs:
(i) a precisely selected coordinate budget and its orthogonal shell,
(ii) an invariant baseline $H_0$ and a cross-block perturbation $V$,
and (iii) control of coupled energy denominators.  Earlier sections
provided exact finite CAR words, charge-selection rules, and conditional
leakage witnesses for spinful perturbations.  Those statements do not
by themselves furnish the spectral decomposition, denominator bounds,
or block invariance needed to insert a particular interacting
Hamiltonian into the Schrieffer--Wolff theorem.

The separately verified finite support results cover the non-resonant
commutator equation and skew-Hermiticity of the constructed generator,
two- and three-level matrix witnesses, the compressed \emph{degree-two
BCH polynomial} identity, the decomposition of shell projections and
their orthogonality and adjoint properties, and a nonzero repeated-mode-one
shell witness. One projector identity is $P+Q=\Pi_{\boldsymbol N,K}$ on the
ambient Fock carrier, whereas the generic compressed BCH theorem assumes
$P+Q=I$ on its whole carrier. Restricting to
$\mathcal B_{\boldsymbol N,K}$ and transferring an actual model's
operators and hypotheses is a separate bridge, not established merely
by placing the two statements next to each other.  These finite support
proofs do not constitute formal verification of the present Chapter~21
manuscript.

The remaining constructive tasks include proving, for any chosen
spinful Hamiltonian, that its finite blocks act on the proposed
carriers; checking charge-sector and energy-margin constraints at
each intermediate application; constructing a valid $H_0$ and its
adapted eigenbasis or another certified Sylvester solver; and proving
any claimed uniform denominator or remainder bounds.  The
coupling-wise condition is conditional, and a vanishing second-order
term is not a substitute for checking it.  Furthermore, a Mott or
spin gap would require an independently defined spectral gap and a
proof of its behavior under relevant size or coupling limits.
A continuum RG flow requires a scale-changing construction and a
comparison between effective theories at successive scales; energy
shell projection at one finite $K$ is not such a construction.
Finally, an isolated finite Hermitian Hamiltonian has a real discrete
spectrum, so scattering lifetimes or dissipative flow cannot be
inferred solely from the matrix identities above.

What has been established here is therefore deliberately narrower
and more precise: a well-defined \emph{finite energy-shell filter}, a
non-resonant \emph{first-order unitary generator}, an exact
\emph{second-order compressed polynomial identity}, and a finite
analytic remainder estimate with explicit hypotheses.  These results
provide the operator-level starting point for a future controlled
effective-Hamiltonian analysis, without prematurely identifying its
outcome with a charge gap or a continuum renormalization-group law.
